\paragraph{An exact three-class consequence of A.159.}
Let
\[
 F=K_1(\zeta_3),\qquad E=K_2(\zeta_3),\qquad
 \sigma(t)=\zeta_3t,\qquad \tau(t)=t,\quad\tau(\zeta_3)=\zeta_3^{-1}.
\tag{A.159.C1}
\]
Then \(E/F\) is cyclic cubic, \(E^{\langle\tau\rangle}=K_2\), and
\(\tau\sigma\tau^{-1}=\sigma^{-1}\).
For \(q=11,59\), the old quadratic place of \(K_1\) splits into two
places \(P_{q,+},P_{q,-}\) of \(F\), exchanged by \(\tau\).
Each is totally ramified in \(E/F\); denote its unique extension by
\(R_{q,+},R_{q,-}\).  The two classes
\[
 \Delta_q=\bigl([R_{q,+}]-[R_{q,-}]\bigr)_{(3)}
                 \in\operatorname{Cl}(E)\{3\},\qquad q=11,59,
\tag{A.159.C2}
\]
generate a subgroup isomorphic to \((\Z/3)^2\).
In particular this is an unconditional class-group lower bound deduced
from the explicit unit certificate, without a computed class-group
presentation.

\begin{proof}
The residue field of the quadratic place is \(\F_{q^2}\), which contains
\(\mu_3\); thus adjoining \(\zeta_3\) splits that place.  At either
resulting place of \(F\), \(v(d^2\alpha)=2\).  Consequently the
extension \(t^3=d^2\alpha\) is totally ramified of degree three.
Writing \(j\) for extension of ideal classes from \(F\) to \(E\),
we have
\[
 3\bigl([R_{q,+}]-[R_{q,-}]\bigr)
     =j\bigl([P_{q,+}]-[P_{q,-}]\bigr).
\tag{A.159.C3}
\]
First work in the exact quotient \(\operatorname{Cl}(E)/j\operatorname{Cl}(F)\).
Suppose a combination with coefficients \(a_{11},a_{59}\in\{0,1,2\}\)
is zero there.  Choose a fractional ideal \(J\) of \(F\) and
\(v\in E^\times\) so that, as an equality of fractional divisors,
\[
 \operatorname{div}(v)=
 \sum_{q=11,59}a_q(R_{q,+}-R_{q,-})-j\operatorname{div}(J).
\tag{A.159.C4}
\]
Here \(j\operatorname{div}(J)\) denotes its extended ideal divisor,
including its ramification multiplicities.  The right side is fixed by
\(\sigma\), so
\[
 u=\frac{\sigma v}{v}\in\mathcal O_E^\times,\qquad
 N_{E/F}u=1,\qquad w=u\tau(u)\in\mathcal O_{K_2}^\times,
 \qquad N_{K_2/K_1}w=1.
\tag{A.159.C5}
\]
The last norm identity follows by applying the cyclic norm and using
\(\tau\sigma\tau^{-1}=\sigma^{-1}\); its restriction to \(K_2\)
is the actual cubic norm to \(K_1\).

At either \(R_{q,\pm}\), a uniformizer is
\(\pi=t^2/q\), since \(v(t)=2\) and \(v(q)=3\).
It satisfies \(\sigma(\pi)/\pi=\zeta_3^2\).
The automorphism \(\sigma\) acts trivially on the residue field.
For an element of valuation \(k\) we therefore have the exact residue
identity
\[
 \overline{\sigma v/v}=\zeta_3^{2k}.
\tag{A.159.C6}
\]
Equation A.159.C4 gives valuations congruent to \(a_q,-a_q\)
modulo three at \(R_{q,+},R_{q,-}\), because every extended ideal
has valuation divisible by three there.  After transporting the second
residue by \(\tau\), which inverts \(\zeta_3\), this yields
\[
 \overline w\bigm|_{R_{q,+}}
   =\zeta_3^{2a_q}\,\tau(\zeta_3^{-2a_q})
   =\zeta_3^{4a_q}=\zeta_3^{a_q}.
\tag{A.159.C7}
\]
The residue field here is also that of the quadratic place of \(K_2\).
Since \((q^2-1)/3\) is prime to three for both \(q=11,59\), the cubic
residue character of \(w\) is trivial precisely when \(a_q=0\).
A.159.8 says that it is trivial for every such relative norm-one unit.
Thus \(a_{11}=a_{59}=0\), proving independence and exact order three
in the displayed quotient.

Finally A.138 gives
\(\operatorname{Cl}(F)\{3\}=0\) for the constant field
\(F=\Q(\zeta_3,\sqrt[3]{3})\).
Hence \(j\operatorname{Cl}(F)\) has order prime to three, and the
quotient map identifies the three-primary parts.  Projecting A.159.C3
to them proves that the two classes in A.159.C2 are killed by three;
the independence already proved shows that they span \((\Z/3)^2\).
\end{proof}

The finite sign and residue checks are retained in
\path{code/two_inert_class_corollary_certificate.py}
and its output.  The proof uses the literal divisor relation in A.159.C4;
it does not identify a relation modulo extended classes with a principal
relation before inserting the fractional ideal \(J\).
